\documentclass[handout]{beamer}
\mode<presentation>
\usepackage{xcolor}
\usepackage[toc]{multitoc}
\usepackage{color}
\usepackage{graphicx}
\graphicspath{{../buku_ajar/buku_ajar_logika_matematika/figures/}}
\usepackage{tikz}
\usetikzlibrary{shapes,arrows,positioning,calc}
\usepackage{listings}
\usepackage{lmodern}
\usepackage{amsmath}
\usetheme[secheader]{Boadilla}
\usecolortheme{whale}
\setbeamertemplate{caption}[numbered]
\setbeamertemplate{bibliography item}[text]
\setbeamertemplate{navigation symbols}{}

\newcommand{\logic}[1]{\textbf{\textit{#1}}}

\theoremstyle{definition}
\newtheorem{definisi}[theorem]{Definisi}
\theoremstyle{example}
\newtheorem{contoh}[theorem]{Contoh}

\renewcommand*{\multicolumntoc}{2}

\setbeamertemplate{footline}
{
  \leavevmode%
  \hbox{%
  \begin{beamercolorbox}[wd=.333333\paperwidth,ht=2.25ex,dp=1ex,center]{author in head/foot}%
    \usebeamerfont{author in head/foot}\insertshortauthor%
  \end{beamercolorbox}%
  \begin{beamercolorbox}[wd=.433333\paperwidth,ht=2.25ex,dp=1ex,center]{title in head/foot}%
    \usebeamerfont{title in head/foot}\insertshorttitle
  \end{beamercolorbox}%
  \begin{beamercolorbox}[wd=.233333\paperwidth,ht=2.25ex,dp=1ex,right]{date in head/foot}%
    \usebeamerfont{date in head/foot}\insertshortdate{}\hspace*{2em} \insertframenumber{} / \inserttotalframenumber\hspace*{2ex} 
  \end{beamercolorbox}}%
  \vskip0pt%
}

\subtitle[Logika Matematika]{Logika Matematika}
\title[Rangkaian Gerbang]{Representasi Fungsi Logika ke Rangkaian Gerbang}
\author{Cecep Suwanda, S.Si., M.Kom.}
\date{Maret 2025}
\institute{Universitas Bale Bandung}

\begin{document}

\maketitle

\section*{Daftar Isi}
\begin{frame}
\frametitle{Daftar Isi}
\tableofcontents
\end{frame}

% ============================================================
% SECTION 11-1: Gerbang Dasar
% ============================================================
\section{Gerbang Logika Dasar}

\begin{frame}
\frametitle{AND, OR, NOT}
\begin{definisi}[AND]
Output 1 \textbf{jika dan hanya jika} semua input 1. $F = A \cdot B$
\end{definisi}
\begin{definisi}[OR]
Output 1 jika \textbf{minimal satu} input 1. $F = A + B$
\end{definisi}
\begin{definisi}[NOT]
Output = komplemen input. $F = A'$
\end{definisi}
Himpunan $\{$AND, OR, NOT$\}$ \textbf{functionally complete}.
\end{frame}

% ============================================================
% SECTION 11-2: Gerbang Turunan
% ============================================================
\section{Gerbang Turunan}

\begin{frame}
\frametitle{NAND, NOR, XOR, XNOR}
\begin{center}
\small
\begin{tabular}{|l|l|l|}
\hline
\textbf{Gerbang} & \textbf{Ekspresi} & \textbf{Fungsi} \\ \hline
NAND & $(AB)' = A'+B'$ & Output 0 hanya jika semua input 1 \\ \hline
NOR & $(A+B)' = A'B'$ & Output 1 hanya jika semua input 0 \\ \hline
XOR & $A \oplus B = A'B+AB'$ & Output 1 jika input berbeda \\ \hline
XNOR & $A \odot B = AB+A'B'$ & Output 1 jika input sama \\ \hline
\end{tabular}
\end{center}
\end{frame}

\begin{frame}
\frametitle{Tabel Kebenaran 2-Input}
\begin{center}
\small
\begin{tabular}{|c|c||c|c|c|c|c|c|}
\hline
$A$ & $B$ & AND & OR & NAND & NOR & XOR & XNOR \\ \hline
0 & 0 & 0 & 0 & 1 & 1 & 0 & 1 \\ \hline
0 & 1 & 0 & 1 & 1 & 0 & 1 & 0 \\ \hline
1 & 0 & 0 & 1 & 1 & 0 & 1 & 0 \\ \hline
1 & 1 & 1 & 1 & 0 & 0 & 0 & 1 \\ \hline
\end{tabular}
\end{center}
\end{frame}

% ============================================================
% SECTION 11-3: Gerbang Universal
% ============================================================
\section{Gerbang Universal}

\begin{frame}
\frametitle{NAND dan NOR: Universal}
\begin{block}{NOT dari NAND}
$A \uparrow A = (AA)' = A'$
\end{block}
\begin{block}{AND dari NAND}
$AB = ((AB)')' = (A \uparrow B) \uparrow (A \uparrow B)$
\end{block}
\begin{block}{OR dari NAND}
$A+B = (A'B')' = (A \uparrow A) \uparrow (B \uparrow B)$
\end{block}
Hanya NAND dan NOR yang universal secara individual.
\end{frame}

\begin{frame}
\frametitle{Konversi ke NAND-Only}
SOP (AND-OR) $\rightarrow$ NAND-NAND:\\
$F = P_1 + P_2 + \ldots = (P_1' \cdot P_2' \cdot \ldots)'$ (De Morgan)\\
Contoh: $F = AB + CD$ $\Rightarrow$ 3 gerbang NAND
\end{frame}

% ============================================================
% SECTION 11-4: Konversi Ekspresi
% ============================================================
\section{Konversi ke Rangkaian}

\begin{frame}
\frametitle{Prosedur Konversi}
\begin{enumerate}
  \item Identifikasi operator terluar (prioritas terendah) = gerbang output
  \item Dekomposisi rekursif: setiap operand dianalisis
  \item Komplemen $\Rightarrow$ gerbang NOT
  \item Variabel dipakai di banyak tempat $\Rightarrow$ fan-out
\end{enumerate}
\end{frame}

\begin{frame}
\frametitle{Contoh: $F = A'B + AB'$ (XOR)}
Operator terluar: OR. Operand kiri $A'B$: AND dengan NOT pada $A$. Operand kanan $AB'$: AND dengan NOT pada $B$.\\
Total: 2 NOT + 2 AND + 1 OR = 5 gerbang.
\end{frame}

% ============================================================
% SECTION 11-5: Half-Adder, Full-Adder
% ============================================================
\section{Rangkaian Aritmetika}

\begin{frame}
\frametitle{Half-Adder}
\begin{definisi}[Half-Adder]
Menjumlahkan 2 bit ($A$, $B$). Output: \textbf{Sum} $S = A \oplus B$, \textbf{Carry} $C_{out} = AB$.
\end{definisi}
\begin{center}
\small
\begin{tabular}{|c|c||c|c|}
\hline
$A$ & $B$ & $S$ & $C_{out}$ \\ \hline
0 & 0 & 0 & 0 \\ \hline
0 & 1 & 1 & 0 \\ \hline
1 & 0 & 1 & 0 \\ \hline
1 & 1 & 0 & 1 \\ \hline
\end{tabular}
\end{center}
1 XOR + 1 AND
\end{frame}

\begin{frame}
\frametitle{Full-Adder dan Ripple-Carry}
\begin{definisi}[Full-Adder]
Menjumlahkan 3 bit: $A$, $B$, $C_{in}$. $S = A \oplus B \oplus C_{in}$, $C_{out} = AB + C_{in}(A \oplus B)$
\end{definisi}
\begin{block}{Ripple-Carry Adder}
$n$ full-adder dirangkai kaskade; carry mengalir dari LSB ke MSB. Delay $\approx O(n)$.
\end{block}
\end{frame}

% ============================================================
% SECTION 11-6: MUX dan Decoder
% ============================================================
\section{Multiplexer dan Decoder}

\begin{frame}
\frametitle{Multiplexer}
\begin{definisi}[MUX]
Memilih satu dari $2^n$ input berdasarkan $n$ sinyal pemilih. MUX $2^n$:1 punya $2^n$ input data, $n$ select, 1 output.
\end{definisi}
MUX 2:1: $Y = S'D_0 + SD_1$\\
MUX 4:1: $Y = S_1'S_0'D_0 + S_1'S_0 D_1 + S_1 S_0' D_2 + S_1 S_0 D_3$
\end{frame}

\begin{frame}
\frametitle{Decoder}
\begin{definisi}[Decoder]
Mengubah kode $n$-bit menjadi $2^n$ output; tepat satu output aktif. Decoder $n$-ke-$2^n$.
\end{definisi}
Setiap output decoder = satu minterm. Fungsi Boolean: decoder + OR (hubungkan output minterm yang relevan).
\end{frame}

% ============================================================
% SECTION 11-7: Analisis Rangkaian
% ============================================================
\section{Analisis Rangkaian}

\begin{frame}
\frametitle{Prosedur Analisis}
\begin{enumerate}
  \item Beri label semua sinyal (input, output, antara)
  \item Telusuri dari input ke output
  \item Substitusi berturut-turut
  \item Sederhanakan; verifikasi tabel kebenaran
\end{enumerate}
\end{frame}

\begin{frame}
\frametitle{Critical Path}
\begin{definisi}[Critical Path]
Jalur sinyal dari input ke output yang melewati gerbang terbanyak. Delay total = jumlah delay gerbang pada jalur kritis. Menentukan frekuensi operasi maksimum.
\end{definisi}
4 tingkat $\times$ $t_d$ = $4 t_d$ delay total.
\end{frame}

% ============================================================
% RANGKUMAN
% ============================================================
\section{Rangkuman}

\begin{frame}
\frametitle{Rangkuman}
\begin{itemize}
  \item AND, OR, NOT: dasar; NAND, NOR: universal
  \item Konversi SOP $\rightarrow$ rangkaian AND-OR; SOP $\rightarrow$ NAND-NAND
  \item Half-adder: $S=A\oplus B$, $C=AB$; Full-adder: + $C_{in}$
  \item MUX: pilih 1 dari banyak; Decoder: kode $\rightarrow$ aktivasi minterm
  \item Critical path menentukan kecepatan
\end{itemize}
\end{frame}

\begin{frame}
\frametitle{Kata Kunci}
\begin{center}
\small
gerbang logika $\bullet$ NAND $\bullet$ NOR $\bullet$ XOR $\bullet$ universal $\bullet$ half-adder $\bullet$ full-adder $\bullet$ ripple-carry $\bullet$ multiplexer $\bullet$ decoder $\bullet$ critical path
\end{center}
\end{frame}

\end{document}
